%-------------------------------------------------------------------------------
% Colter Snyder -- Cover letter, Google Software Engineer III, Infrastructure,
% Platforms Infrastructure Engineering
% Build: latexmk -pdf -outdir=build ColterSnyder_CoverLetter_Google_SWE3_Infra.tex
%-------------------------------------------------------------------------------
\documentclass[letterpaper,11pt]{article}

\usepackage[margin=0.8in,top=0.6in,bottom=0.6in]{geometry}
\usepackage[T1]{fontenc}
\usepackage{lmodern}
\usepackage{xcolor}
\usepackage{fontawesome5}
\usepackage[hidelinks]{hyperref}

\definecolor{accent}{HTML}{1F5FA8}
\pagestyle{empty}
\setlength{\parindent}{0pt}
\setlength{\parskip}{8pt}
\pdfgentounicode=1

\begin{document}

%--- Header (matches resume) ---------------------------------------------------
\begin{center}
  {\Huge\bfseries Colter Snyder}\par\vspace{4pt}
  \footnotesize
  Lakewood, CO \enspace$\cdot$\enspace
  \faPhone*\ (720) 626-2558 \enspace$\cdot$\enspace
  \faEnvelope\ \href{mailto:coltersnyder@gmail.com}{coltersnyder@gmail.com} \enspace$\cdot$\enspace
  \mbox{\faLinkedin\ \href{https://www.linkedin.com/in/colter-snyder-a36854183/}{in/colter-snyder-a36854183}}
\end{center}
{\color{accent}\hrule height 0.6pt}
\vspace{6pt}

October 23, 2026

Hiring Team\\
Google

\textbf{Re: Software Engineer III, Infrastructure, Platforms Infrastructure Engineering}

Dear Hiring Team,

The infrastructure behind Google's AI and compute platforms spans everything from service code to
TPUs, and that full stack is exactly the range I enjoy working across. I am a software engineer at
Innovative Scientific Solutions, where I build back-end calculation services for simulation platforms
used by defense agencies. I hold an M.S.\ in Computer Science from the Colorado School of Mines, and I
would like to bring my experience with services, performance, and hardware to Platforms
Infrastructure Engineering.

At ISSI I build Rust and Python services that communicate over WebSockets, store data in MongoDB, and
deploy in Docker containers on Linux. When the team needed to understand how those services performed,
I spearheaded an internal benchmarking tool in Python, first as a command-line application and then
with a GUI front end. It runs any combination of our 13 effect calculations and charts the results,
and it exposed a tenfold slowdown in our interface to an external library, with graphs clear enough to
take straight to leadership. It also saved what would likely have been weeks of manual testing with
hand-written JSON inputs. Beyond performance work, I port legacy physics
models from Fortran and Python to Rust while preserving their numerical behavior, and I lead development of
individual features from design documents through automated tests. I review teammates' merge
requests regularly, since every change requires a separate reviewer and merger.

My interest in infrastructure goes below the operating system. In my own time I design FPGA logic in
Verilog, program in C, and write 6502 and Z80 assembly, which has given me a working sense of how
software meets compute hardware. I have also co-taught Data Structures and Algorithms in C++ to 120
students, and diagnosing a 2,000-test suite during my Tyler Technologies internship, where I restored
about 1,500 tests to passing, taught me to track problems to their source rather than their symptoms.

I am eager to grow into larger-scale distributed systems, and I can think of no better place to do it
than the team building Google's computing foundation. I live in Lakewood and would be glad to work from Google's Boulder office or remotely. Thank you for your time and consideration.

Sincerely,

Colter Snyder

\end{document}
